\chapter{Introduction}
This dissertation deals with different ways of detecting semantic relations in text corpora. Its focus is more on methodology than on analytical results. Some important restrictions are necessitated by the precise research question: In a corpus of Latin texts, how can we automatically extract semantic relations that are relevant for language learning? I will concentrate only on a subset of the existing phenomena, namely hypernymy and hyponymy. This selection is motivated by the objective of utilizing the obtained insights as the basis for exercises in the process of learning historical languages. Thus, the detection of said relations is not an end in itself, but is to be viewed in the context of how humans learn languages. 

\section{Semantic relations in language learning}
Specifically, recent studies suggest that an improved integration of hyponymy into the learning process can lead to a higher level of language skills \parencite[196]{taslimExperimentalStudyTeaching2014}. Roughly the same procedure is applicable also in the context of artificial intelligence, which highlights the general importance of semantic relations in any kind of language learning effort \parencite[1306]{carlsonArchitectureNeverendingLanguage2010}. The automated extraction of semantic relations from texts can be used for constructing ontologies in specific domains, which is important for language learning because many textbooks are heavily structured according to a range different topics \parencite[192f.]{punuruLearningNontaxonomicalSemantic2012}. 

Since there is a parallel between semantic relations (or ontologies in general) and the way that humans think, it seems reasonable to exploit that connection for learning purposes. This is particularly important for advanced learners' lexical acquisition \parencite[56]{crossleyDevelopmentSemanticRelations2010}. Now, one might be tempted to conclude that predefined lists of basic vocabulary and more advanced semantic relations would be a sufficient help to language teachers. However, static materials und human supervision will probably fail to account for the entire processual complexity of lexical acquisition \parencite[71]{crossleyDevelopmentSemanticRelations2010}. The involvement of artificial intelligence in the modeling of human knowledge about vocabulary faces various problems, including the progression of learning over time \parencite[70]{crossleyDevelopmentSemanticRelations2010}. 

In language learning contexts, this problem extends also to the morphological domain, where it can prove difficult for students to exploit their knowledge about affixation for the acquisition of semantic relations \parencite[248]{gardnerValidatingConstructWord2007}. However, "inflected and derived forms" are exactly what we encounter in the curricular targets for Latin courses, i.e. literary texts. Most textbooks take care of the problem by using hand-crafted Latin texts instead of authentic ones. This leads to a highly problematic transition from textbook to literature, which is already famous among German teachers of Latin \parencite[115]{schibelZurAneignungLateinischer2013}. The challenge, then, is to not only extract semantic relations from a given corpus, but also to group them according to general principles of progression that are tightly coupled with a domain-specific theory of learning.